\section{Online Experiments}

The term crowd-sourcing emerged analogously to the term outsourcing, which describes the process of relocating production locations to other countries mostly with lower production costs. The initial idea of crowd-sourcing was to recruit groups of internet users to complete human intelligence tasks (HITs), instead of hiring temporary employees. These tasks can vary from commercial applications such as attaching meaningful tags to numerous products of a new online shop to cleaning up training data for machine learning or spectacular database investigations \parencite{CrowdsourceML,TurkRescue}. Specifically, these are tasks that are too complicated to be solved by computers and need human intelligence but do not require skilled personnel \parencite{BerendViabCrowd2011}. 
%Crowd sourcing can be defined as “the intentional mobilization for commercial exploitation of creative ideas and other forms of work performed by consumers” \parencite{def-crowdsource}. 
Importantly, crowd source workers are paid, recruited from any geographic location and hired only for single tasks. 

In recent years it has been discovered that especially science can benefit from this new approach. Online experiments offer the opportunity to collect very simple, flexible and fast data from a large pool of participants \parencite{BerendViabCrowd2011, truellResponse2002}. Reduced costs, easier logistics, and lower investment of physical resources are other advantages of online experiments \parencite{KrautOnlExpCost}. But the most significant benefit offered by conducting experiments online is the potential to surmount the prevalent reliance on a participant pool characterized as WEIRD. Online participant pools are found to be older as well as more ethnically and educationally diverse and thereby may help to enhance the diversity and representativeness of research samples \parencite{BerendViabCrowd2011, henrichWeirdest2010}. Earlier studies have also shown that the pools of online providers such as Mechanical Turk are more representative for the average working adult population within a country \parencite{BerendViabCrowd2011}. This is crucial as \textcite{WardUndergradsEmployed1993} has indicated even significant differences between undergraduate and employed American adults.
% At the same time data are still similar to the ones obtained from undergraduate students and are reported to have comparable quality to in lab experiments \parencite{BerendViabCrowd2011, sprouseValidAMazTurk2010}.
% providers

%MTurk
The first relevant provider that offered the opportunity to recruit participants online was amazon mechanical turk (MTurk, \parencite{mturk}). In 2015 MTurk generously claimed to have 500,000 online workers available, which was corrected down to 7,300 actually recruitable participants, which is still a large participant pool \parencite{stewartAverage2015}. Although the platform was not intended to host research surveys, in 2015 already 1,200 studies were conducted online. Some of these studies replicated results from older lab studies and therefore proved consistency across the procedures \parencite{CrumpEval2013, AmirEcoGames2012, hortonOnline2011}.

%Prolific
However, several problems related to the pool management of MTurk were discovered. This got addressed by Prolific \parencite{prolific} which is currently the predominant recruiter used for online experiments providing access to a pool of currently 120,000 participants from 35 countries %(almost all from the OECD-list).
In contrast, Prolific was specifically designed for research surveys and informs participants accordingly. Again comparable results show consistency with Mturk \parencite{peerTurk2017}. Prolific ensures fair payment, and gives incentives for participants to provide qualitative answers. The participant pool is even more geological and ethical diverse and the website design allows to target participants based on demographic criteria. Each participant is strictly identified to ensure that each person can only participate once and to exclude bot participants. Conversation between experimenter and participants is monitored by Prolific to prohibit demand effects and exploitation \parencite{NeckaMturkProblematic2016}.

% Lucid
The geographical reach of Prolific predominantly encompasses countries listed by the OECD, which implies that surveys carried out through this platform are largely confined to industrialized nations. Conversely, Lucid \parencite{lucid} serves as an alternative marketplace, covering a broader spectrum of 112 countries and 73 languages that include both industrialized and non-industrialized regions, offering a more diverse participant pool. Importantly, participants on Lucid are approached using their own native language whereas all communication on Prolific is in English.
This opens up new possibilities to reach an even broader pool of participants around the world. To our knowledge, this is the first research survey to use Lucid as a recruiting platform and to take advantage of some of this potential.
The task was not only to identify the recruiting mechanisms of this platform, but also to evaluate its consistency compared to established platforms such as Prolific.




% Lucid is not tailred towards research

% Both recruiters allow the use of individual software to acustom the surveys






% Because of poor participants moderation the quality of this provider declined \parencite{Kennedy_Clifford_Burleigh_Waggoner_Jewell_Winter_2020}.
% Since then a commonly used provider is Prolific. 


% Loss of nativity:
% - participation in hundreds of studies may have a practice effect.
% - Dicussion board organizes exchange between MTurk workers
% - Conversation can be monitored but not controlled 
% No controlled in lab experiment conditions
% - Multitasking (watching TV or listen to music Chandler et al. (2014)Necka et al. (2016)
% - No substantial differences between groups found Necka et al. (2016)
% - Caution for experiments if design requires stable conditions across participants
% Payment:
% - Ethical problems in the case of exploiting participants by paying to less as well as overpaying and therefore luring them into participation of experiment they would not do for lower wage. Shank, 2016; Crump et al., 2013; Mason and Suri, 2012)
% - Although data quality seems to be independent from reward size it might still have a arbitrary influence on the data quality (Buhrmester et al., 2011)
% Participant IDs:
% - Tracking of participants that use multiple accounts and therefore participate more than once in the survey
% - (Chandler et al., 2014).
% Demand effect:
% - To maintain high data quality the experimenter is able to give feedback or even reject the payment of participants according to the answers of a participant. The feedback results in an acceptance score assigned to each participant wich can then lead to exclusion to certain surveys. 
% - To avoid low ratings participants might give answers according to how they expect the results of the survey are.
% - Necka et al. (2016)


% Prolific:
% 1.2 dedicated research pool:
% - explicetly recruits participants for research and informs them accordingly.
% - Study have replicated results of experiments conducted on Mturk
% - Reported a geographical and ethnical more diverse pool

% 2. Subject pool management:
% - Ensure fair payment by setting a minimum fixed payment of 5 GBP per hour.
% - In fact this might be problematic in case of developing countries as it is considered unethical to offer excessive incentives to lure participants into studies
% - Transparent information of the duration of a experiment
% - Communication between provider and participant is mediated by Prolific. 
% - Rejection of low quality data have to be justified by the experimenter and participants can raise an objection in case.
% - This lowers the danger of demand effects that the participant answers as he thinks he is supposed to answer.
% - In addition participant always can terminate or return their submission without fear of negative reputation 
% - Very good participants can even win prizes
% - H

% 3. Pre-Screening:
% - Experimenters can exclude participants based on pre screening questions. 
% - To mitigate the risk of response bias, where participants might answer based on perceived expectations, basic demographic information is collected at the account creation stage. Furthermore, these responses are stored and made available for subsequent experiments, ensuring consistency and reducing the need for repetitive data collection.

% 4. Exclusion of individual participants
% - Prolific employs a system where each participant is assigned a unique participant ID and monitors their activity on the platform. This approach enables the prevention of individuals from participating in the same experiment more than once, ensuring the integrity of experimental data.

